\documentclass[11pt,a4paper]{article}

%  XeLaTeX
\usepackage{fontspec}

\usepackage{algorithm}
\usepackage{algorithmic}
\usepackage{amsmath,amssymb,amsthm}
\usepackage{bm}
\usepackage[useregional]{datetime2}
\usepackage{enumitem}
\usepackage{booktabs}
\usepackage{float}
\usepackage{graphicx}
\usepackage{tabularray}
\usepackage{framed}
\usepackage[a4paper, margin=2.5cm]{geometry}
\usepackage{eso-pic}
\usepackage{mathtools}
\usepackage{tikz}
\usetikzlibrary{calc}
\usepackage{xcolor}

\usepackage[backend=biber, style=apa, sorting=nyt]{biblatex}
\addbibresource{D:/github/ENEL445/report/references.bib}
\graphicspath{{D:/github/ENEL445/figs/}}

\usepackage[colorlinks=true, linkcolor=blue, citecolor=blue, urlcolor=blue]{hyperref}
\hypersetup{
  pdfauthor   = {David Ewing},
  pdfsubject  = {\jobname.tex},
  pdfkeywords = {source: \jobname.tex}
}

% Running margin label
\newcommand{\mymarginlabel}{\textcopyright\ 2026 David Ewing dew59@uclive.ac.nz | \DTMnow\ | \jobname.tex}
\AddToShipoutPictureBG{%
  \begin{tikzpicture}[remember picture,overlay]
    \node[rotate=90, anchor=north]
      at ($(current page.north west)+(1.4cm,-13cm)$)
      {\scriptsize\ttfamily \mymarginlabel};
  \end{tikzpicture}%
}

% Custom commands
\newcommand{\E}{\mathbb{E}}
\newcommand{\R}{\mathbb{R}}
\newcommand{\N}{\mathcal{N}}
\newcommand{\ELBO}{\text{ELBO}}
\newcommand{\tr}{\text{tr}}
\newcommand{\sigmoid}{\sigma}
\setlength{\parindent}{0pt}

\title{
Posterior Variance Collapse in Variational Inference:\\
Diagnosis and Laplace Correction Across Three Bayesian Regression Models
}
\author{
  David Ewing (82171165)\\
  ENEL445 --- Applied Engineering Optimisation
}
\date{15 May 2026}

\begin{document}

\maketitle

\begin{abstract}
\setlength{\parindent}{0pt}
\setlength{\parskip}{0.7em}

Mean-field variational inference turns intractable posterior computation into an
optimisation problem by maximising the Evidence Lower Bound (ELBO), but the
mean-field factorisation systematically underestimates posterior uncertainty ---
a problem known as \emph{posterior variance collapse}.
This report diagnoses that phenomenon across three Bayesian regression models and
proposes a Laplace approximation correction to recover more faithful credible intervals.

Five case studies were conducted (linear, quadratic, logistic, hierarchical linear,
and hierarchical logistic regression).
The quadratic (CS2) and logistic (CS3) cases are excluded from the correction
analysis because their collapse is negligible
($\sigma_{\mathrm{VB}}/\sigma_{\mathrm{Gibbs}} > 0.97$ for all parameters), so the
focus falls on \textbf{CS1} (Bayesian linear regression),
\textbf{CS4} (hierarchical linear), and \textbf{CS5} (hierarchical logistic).

In each retained case, four optimisation methods --- CAVI, gradient ascent, Newton,
and BFGS --- are used to maximise the ELBO.
Once CAVI has converged, its posterior mean $\hat{\bm{\mu}}$ is accepted as the best
available point estimate, and the covariance is corrected by inverting the Hessian
of the joint log-posterior evaluated at that point:
\begin{equation*}
  \hat{\bm{\Sigma}}_{\mathrm{corr}}
  = \bigl(-\nabla^2 \log p(\bm{\theta}\mid\bm{y})\bigr)^{-1}\!\Big|_{\bm{\theta}=\hat{\bm{\mu}}}.
\end{equation*}
The Hessian is computed by central finite differences at a single point; CAVI is
\emph{not} re-run, so the additional cost is negligible.
Corrected credible intervals are then compared against the Gibbs gold standard.

The correction is effective in both hierarchical settings.
For CS4 (hierarchical linear), the intercept SD ratio rises from $0.130$ to $0.764$
while the slope already reaches $0.998$.
For CS5 (hierarchical logistic), both coefficients move from $\approx 0.46$ to
$0.82$--$0.88$.
In CS1 the collapse is mild to begin with ($\approx 0.98$) and the correction changes
nothing, which is reassuring: it confirms that the problem is structural rather than
a numerical artefact.
\end{abstract}

\newpage
\tableofcontents
\newpage

%─────────────────────────────────────────────────────────────────────────────
\section*{Introduction}
\addcontentsline{toc}{section}{Introduction}
%─────────────────────────────────────────────────────────────────────────────

\subsection*{Bayesian inference and the intractability problem}

Bayesian regression places a prior distribution on parameters $\bm{\theta}$ and updates
it via the likelihood to obtain the posterior:
\begin{equation}
  p(\bm{\theta}\mid\bm{y})
  = \frac{p(\bm{y}\mid\bm{\theta})\,p(\bm{\theta})}
         {\displaystyle\int p(\bm{y}\mid\bm{\theta})\,p(\bm{\theta})\,\mathrm{d}\bm{\theta}}.
  \label{eq:posterior}
\end{equation}
The denominator of \eqref{eq:posterior} is analytically tractable only for conjugate
likelihood--prior pairs.
For all three models studied here the integral is intractable, though for different reasons:
hierarchical models introduce high-dimensional latent variables, and the logistic
likelihood has no conjugate prior.

\medskip
The three models studied are shown below in parallel, with each difference
highlighted:

\begin{align}
  \text{CS1 (linear):}\quad
  & y_i \mid \bm{\beta}, \tau_e
    \;\sim\; \mathcal{N}\!\left(
      \underbrace{\bm{x}_i^T\bm{\beta}}_{\text{fixed effects}},\;
      \tau_e^{-1}
    \right),
    \quad
    p(\bm{\beta},\tau_e)\propto 1\;\text{(flat)},
    \nonumber\\[6pt]
  \text{CS4 (hier.\ linear):}\quad
  & y_{ij}\mid\bm{\beta},u_j,\tau_e
    \;\sim\; \mathcal{N}\!\left(
      \underbrace{\bm{x}_{ij}^T\bm{\beta}+u_j}_{\text{fixed + random}},\;
      \tau_e^{-1}
    \right),\;
    u_j\!\sim\!\mathcal{N}(0,\tau_u^{-1}),\;
    \bm{\beta}\!\sim\!\mathcal{N}(\bm{0},\lambda^{-1}\bm{I}),
    \nonumber\\[6pt]
  \text{CS5 (hier.\ logistic):}\quad
  & y_{ij}\mid\bm{\beta},u_j
    \;\sim\;
    \underbrace{\mathrm{Bernoulli}}_{\text{non-Gaussian}}\!\!\left(
      \sigma\!\!\left(
        \underbrace{\bm{x}_{ij}^T\bm{\beta}+u_j}_{\text{fixed + random}}
      \right)
    \right),\;
    u_j\!\sim\!\mathcal{N}(0,\tau_u^{-1}),\;
    \bm{\beta}\!\sim\!\mathcal{N}(\bm{0},\lambda^{-1}\bm{I}).
    \nonumber
\end{align}

\noindent
In all three models $\bm{\beta}\in\R^2$ with
$\bm{x}_i = (1, x_i)^T$,
and the group index $j\in\{1,\ldots,5\}$ applies to CS4 and CS5.
CS4 and CS5 share the same random-effect structure; they differ in whether the
observation model is Gaussian (CS4, with additional noise precision $\tau_e$) or
Bernoulli-logistic (CS5, with no explicit noise variance).

\subsection*{Scope of this report}

A total of five case studies were conducted, covering linear, quadratic, logistic,
hierarchical linear, and hierarchical logistic regression.
\textbf{Case Studies 2 (quadratic) and 3 (logistic) are excluded from the correction
analysis} because their posterior SD ratios $\sigma_{\mathrm{VB}}/\sigma_{\mathrm{Gibbs}}$
are greater than $0.97$ for all parameters --- collapse is negligible and no correction
is warranted.
The analysis focuses on CS1, CS4, and CS5, where collapse ranges from mild to severe.

\subsection*{Variational Inference as optimisation}

Variational Inference (VI) replaces \eqref{eq:posterior} with the optimisation problem
\begin{equation}
  \max_{q\in\mathcal{Q}}\,\mathcal{L}(q)
  = \E_{q}\!\left[\log\frac{p(\bm{y},\bm{\theta})}{q(\bm{\theta})}\right]
  = \E_q[\log p(\bm{y}\mid\bm{\theta})] - \mathrm{KL}(q\,\|\,p),
  \label{eq:elbo}
\end{equation}
where $\mathcal{L}(q)$ is the Evidence Lower Bound (ELBO).
Maximising \eqref{eq:elbo} is equivalent to minimising the KL divergence between
the approximation $q$ and the true posterior.
In the mean-field approximation the posterior is factorised as
$q(\bm{\theta}) = q(\bm{\beta})\,q(\tau_e)\,\prod_j q(u_j)\,q(\tau_u)$,
and each factor is optimised in turn (CAVI) or jointly via gradient-based methods.

Four optimisation methods are applied to the ELBO in each case study:
CAVI, gradient ascent with Armijo backtracking, Newton's method with finite-difference
Hessian, and BFGS with Wolfe conditions.
Reference posteriors are obtained from Gibbs samplers (conjugate for linear models;
P\'{o}lya--Gamma \parencite{Polson2013} for logistic models).

\subsection*{The variance-collapse problem and the Laplace correction}

A known failure mode of mean-field VI is \emph{posterior variance collapse}: the
optimised approximation $q^\star$ systematically underestimates posterior uncertainty
because the mean-field factorisation eliminates posterior correlations and reduces
effective parameter uncertainty.
Concretely, the CAVI update for the $\bm{\beta}$ covariance in CS1 is
\begin{equation*}
  \bm{\Sigma}_\beta = \bigl(\E[\tau_e]\,\bm{X}^T\bm{X} + \lambda\bm{I}\bigr)^{-1},
\end{equation*}
which treats $\tau_e$ as fixed at its variational mean rather than integrating over its
uncertainty; this systematically tightens the covariance.
In hierarchical models the collapse is amplified because the variational distribution
over group effects $u_j$ is also factorised from $\bm{\beta}$.

The remedy proposed here is the \textbf{Laplace approximation} centred at the CAVI
solution.
The CAVI mean $\hat{\bm{\mu}}$ is accepted as the best available point estimate; the
full joint log-posterior $\log p(\bm{\theta}\mid\bm{y})$ is then written in unconstrained
coordinates $(\bm{\beta},\,\mathbf{u},\,\log\tau_e,\,\log\tau_u)$ and its Hessian
$H = \nabla^2\log p\,\big|_{\hat{\bm{\mu}}}$ is evaluated by central finite differences
at that single point.
The corrected covariance is $\hat{\bm{\Sigma}}_{\mathrm{corr}} = (-H)^{-1}$.
CAVI is \emph{not} re-run: only the log-posterior curvature is needed, not the ELBO.
The gradient-based methods (Newton, BFGS, gradient ascent) serve to cross-check the
Hessian and confirm that $\hat{\bm{\mu}}$ sits near the log-posterior mode.

\begin{description}
  \item[Case Study 1 --- Linear Regression ($n=50$, $p=2$):]
    Serves as the baseline.
    With CAVI SD ratios of $\approx 0.98$, collapse is mild and the Laplace correction
    leaves the intervals essentially unchanged --- confirming that the problem does not
    arise when the posterior is near-Gaussian.
  \item[Case Study 4 --- Hierarchical Linear ($N=150$, $J=5$, $p=2$):]
    Shows severe asymmetric collapse: the intercept $\beta_0$ has an SD ratio of only
    $0.130$ while the slope is barely affected ($0.987$), because $\beta_0$ is entangled
    with the group random effects in a way that $\beta_1$ is not.
    The Laplace correction lifts $\beta_0$ to $0.764$ and $\beta_1$ to $0.998$.
  \item[Case Study 5 --- Hierarchical Logistic ($N=150$, $J=5$, $p=2$):]
    Exhibits symmetric collapse, with both coefficients stuck near $0.46$ --- a pattern
    that the Laplace correction recovers to $0.82$--$0.88$.
\end{description}

%─────────────────────────────────────────────────────────────────────────────
\section{Common Methodology}
\label{sec:methodology}
%─────────────────────────────────────────────────────────────────────────────

\subsection{Variational Inference as Optimisation}

The mean-field variational approximation factorises the joint posterior as
$q(\bm{\beta}, \tau_e) = q(\bm{\beta})\,q(\tau_e)$ for the Gaussian regression cases,
or $q(\bm{\beta}) = \N(\bm{\mu}, \bm{\Sigma})$ for the logistic case.
Maximising the ELBO over the variational parameters $\bm{\phi}$ is equivalent to
minimising the KL divergence $\mathrm{KL}(q\|p)$ between the approximation and the
true posterior.

\subsection{Constraint Handling: Cholesky Reparameterisation}

The covariance matrix $\bm{\Sigma}$ must remain symmetric positive definite at every
iteration.
The set of symmetric positive definite matrices has a curved boundary: a standard
Euclidean gradient step can cross outside it, producing an invalid covariance.
The solution is to parameterise over the Cholesky factor $\bm{L}$, where
$\bm{\Sigma} = \bm{L}\bm{L}^T$ and $\bm{L}$ is lower-triangular with positive diagonal
entries.
The diagonal entries are reparameterised as $L_{ii} = e^{\tilde{\ell}_{ii}}$ to allow
unconstrained optimisation.
This converts the constraint into a smooth bijection and ensures every gradient step
produces a valid covariance.

\subsection{Optimisation Methods}

All four methods share the same objective (the ELBO), the same reparameterisation,
and the same convergence criteria.

\begin{description}
  \item[CAVI:] Coordinate Ascent Variational Inference.
    Closed-form coordinate updates derived from the exponential family structure.
    Each update exactly maximises the ELBO over one block of parameters while holding
    the others fixed.
    Most efficient method; requires conjugate structure.
  \item[Gradient ascent:] First-order method.
    Search direction is the gradient $\bm{d}^{(t)} = \nabla\mathcal{L}(\bm{\phi}^{(t)})$.
    Step size determined by Armijo backtracking (sufficient increase condition).
  \item[Newton's method:] Second-order method.
    Search direction is $\bm{d}^{(t)} = -\bm{H}^{-1}\nabla\mathcal{L}$,
    where $\bm{H}$ is the Hessian approximated by finite differences (25 evaluations
    per iteration).
    Fast convergence near the optimum; expensive per iteration.
  \item[BFGS:] Quasi-Newton method.
    Builds a curvature approximation from gradient differences, requiring only
    gradient evaluations.
    Step size determined by Wolfe conditions to maintain positive definiteness of
    the approximate inverse Hessian ($\bm{y}_t^T\bm{s}_t > 0$).
    Best practical compromise between convergence speed and cost.
\end{description}

For all line-search methods, the search direction must be an ascent direction:
$\nabla\mathcal{L}(\bm{\phi}^{(t)})^T\bm{d}^{(t)} > 0$.

\subsection{Convergence Criteria}

Iteration terminates when any of the following conditions is satisfied:
\begin{enumerate}[label=(\roman*)]
  \item Gradient norm: $\lVert\nabla\mathcal{L}\rVert < \varepsilon_g$.
  \item Relative ELBO change: $|\Delta\mathcal{L}|/|\mathcal{L}| < \varepsilon_r$.
  \item Parameter change: $\lVert\Delta\bm{\phi}\rVert < \varepsilon_\phi$.
  \item Maximum iterations: $t \geq t_{\max}$.
\end{enumerate}
CAVI additionally monitors the maximum relative change in variational parameters.

\subsection{Reference Samplers}

Two types of Gibbs sampler are used as gold-standard references.

\paragraph{Conjugate Gibbs (linear and quadratic cases).}
The Normal--Gamma conjugate structure of the Bayesian regression model (Gaussian
likelihood with Gaussian prior on $\bm{\beta}$ and Gamma prior on precision $\tau_e$)
admits closed-form Gibbs updates:
\begin{align}
  \bm{\beta} \mid \tau_e, \bm{y} &\sim \N(\bm{\mu}_n, \bm{\Sigma}_n), \\
  \tau_e \mid \bm{\beta}, \bm{y} &\sim \mathrm{Gamma}(a_n, b_n),
\end{align}
where $\bm{\Sigma}_n = (\tau_e\bm{X}^T\bm{X} + \lambda_\beta\bm{I})^{-1}$ and
$\bm{\mu}_n = \tau_e\bm{\Sigma}_n\bm{X}^T\bm{y}$.

\paragraph{P\'{o}lya--Gamma Gibbs (logistic case).}
The P\'{o}lya--Gamma (PG) sampler of \textcite{Polson2013} augments the logistic model
with latent variables $\omega_i \sim \mathrm{PG}(1, x_i^T\bm{\beta})$, rendering
$\bm{\beta}$ conditionally Gaussian:
\begin{align}
  \bm{\beta} \mid \bm{\omega}, \bm{y}
    &\sim \N\!\left(\bm{\Sigma}_\omega(\bm{X}^T\bm{\kappa}),\; \bm{\Sigma}_\omega\right),
  \quad \bm{\Sigma}_\omega = (\bm{X}^T\bm{\Omega}\bm{X} + \bm{\Sigma}_0^{-1})^{-1},
\end{align}
where $\bm{\Omega} = \mathrm{diag}(\omega_i)$ and $\kappa_i = y_i - 1/2$.
The PG draw requires $T$ Gamma variates per latent variable; the truncation $T=50$
is used throughout.

\subsection{Common Notation}

\begin{center}
\begin{tabular}{@{}ll@{}}
\toprule
Symbol & Meaning \\
\midrule
\multicolumn{2}{@{}l}{\textbf{Bayesian model}} \\
$\bm{X}$ & Design matrix ($n \times p$), rows $x_i^T$ \\
$\bm{y}$ & Response vector \\
$\bm{\beta}$ & Regression coefficient vector \\
$\tau_e$ & Observation-noise precision (Gaussian cases) \\
$n,\,p$ & Number of observations; number of predictors \\
$\lambda_\beta$ & Prior precision on $\bm{\beta}$ \\
\midrule
\multicolumn{2}{@{}l}{\textbf{Variational approximation}} \\
$q(\bm{\beta})$ & Variational Gaussian approximation \\
$\bm{\mu},\,\bm{\Sigma}$ & Mean and covariance of $q(\bm{\beta})$ \\
$\bm{L}$ & Cholesky factor, $\bm{\Sigma}=\bm{L}\bm{L}^T$ \\
$\mathcal{L}(\bm{\phi})$ & Evidence Lower Bound (ELBO) \\
\midrule
\multicolumn{2}{@{}l}{\textbf{Optimisation}} \\
$\bm{d}^{(t)}$ & Search direction at iteration $t$ \\
$\alpha_t$ & Step size at iteration $t$ \\
$\bm{s}_t,\,\bm{y}_t$ & Parameter-change and gradient-difference vectors (BFGS) \\
\midrule
\multicolumn{2}{@{}l}{\textbf{Logistic-specific}} \\
$\sigma(\cdot)$ & Logistic sigmoid, $\sigma(\psi)=(1+e^{-\psi})^{-1}$ \\
$\xi_i$ & Jaakkola--Jordan tightness parameter \\
$\omega_i$ & P\'{o}lya--Gamma latent variable \\
$\kappa_i$ & Pseudo-response $\kappa_i = y_i - 1/2$ \\
$\mathrm{ESS}$ & Effective sample size \\
\midrule
\multicolumn{2}{@{}l}{\textbf{Lecture correspondence (Prof.\ R.\,S.\ Smith, ENEL445 2026)}} \\
$J(\theta)$ & Cost function $= \mathcal{L}(\bm{\phi})$ (ELBO) \\
$\Phi$ & Regressor matrix $= \bm{X}$ \\
\bottomrule
\end{tabular}
\end{center}

%─────────────────────────────────────────────────────────────────────────────
\section{Case Study 1: Bayesian Linear Regression}
\label{sec:linear}
%─────────────────────────────────────────────────────────────────────────────

\subsection{Model}

The Bayesian linear regression model is
\begin{align}
  \bm{y} &\sim \N(\bm{X}\bm{\beta},\; \tau_e^{-1}\bm{I}_n), \\
  \bm{\beta} &\sim \N(\bm{0},\; \lambda_\beta^{-1}\bm{I}_p), \quad \lambda_\beta = 0.1, \\
  \tau_e &\sim \mathrm{Gamma}(\alpha_e, \gamma_e).
\end{align}
The mean-field variational approximation is
$q(\bm{\beta},\tau_e) = \N(\bm{\mu}_\beta,\bm{\Sigma}_\beta)\cdot\mathrm{Gamma}(a_e,b_e)$.

The test case uses $n=50$ observations with $x_i \sim \mathrm{Uniform}(-2,2)$ and
true parameters $\bm{\beta}_{\mathrm{true}} = (1.0,\; 2.0)^T$ ($p=2$).
Data are generated with noise standard deviation $\sigma=0.5$.

\subsection{ELBO}

Under the Normal--Gamma variational family, the ELBO is
\begin{align}
  \mathcal{L}(\bm{\phi})
  &= \E_q[\log p(\bm{y}\mid\bm{\beta},\tau_e)]
   + \E_q[\log p(\bm{\beta})]
   + \E_q[\log p(\tau_e)]
   - \E_q[\log q(\bm{\beta})]
   - \E_q[\log q(\tau_e)].
\end{align}
Each term has a closed-form expression:
\begin{align}
  \E_q[\log p(\bm{y}\mid\bm{\beta},\tau_e)]
  &= \frac{n}{2}\bigl(\psi(a_e)-\log b_e\bigr)
    - \frac{a_e}{2b_e}\bigl(
        \lVert\bm{y}-\bm{X}\bm{\mu}_\beta\rVert^2
        + \tr(\bm{X}\bm{\Sigma}_\beta\bm{X}^T)
      \bigr)
    - \frac{n}{2}\log(2\pi),\\
  \mathrm{KL}(q(\bm{\beta})\|p(\bm{\beta}))
  &= \tfrac{1}{2}\bigl[
      \lambda_\beta\,\tr(\bm{\Sigma}_\beta)
      + \lambda_\beta\lVert\bm{\mu}_\beta\rVert^2
      - p
      - \log\det(\lambda_\beta\bm{\Sigma}_\beta)
    \bigr],
\end{align}
with analogous terms for the Gamma factor.
Full derivations are in the individual linear report (Appendix~\ref{app:linear-ref}).

\subsection{Results}

\begin{figure}[H]
  \centering
  \includegraphics[width=0.85\textwidth]{linear_data_scatter.png}
  \caption{Synthetic data for Case Study~1: $n=50$ observations, $x_i\sim\mathrm{Uniform}(-2,2)$,
           true $\bm{\beta}=(1.0,2.0)$, noise $\sigma=0.5$.}
  \label{fig:linear-data}
\end{figure}

\begin{figure}[H]
  \centering
  \includegraphics[width=0.85\textwidth]{linear_elbo_convergence.png}
  \caption{ELBO convergence for all four optimisation methods on the linear test case.
           CAVI converges in approximately 50 iterations; gradient ascent requires 2000.
           BFGS matches CAVI's final ELBO in fewer than 50 iterations.}
  \label{fig:linear-elbo}
\end{figure}

\begin{figure}[H]
  \centering
  \includegraphics[width=0.85\textwidth]{linear_vb_gibbs_beta0.png}
  \caption{Posterior comparison for $\beta_0$: variational approximation
           versus conjugate Gibbs reference density (linear case).
           Mild variance collapse: the VB approximation is slightly overconfident.}
  \label{fig:linear-posterior-beta0}
\end{figure}

\begin{figure}[H]
  \centering
  \includegraphics[width=0.85\textwidth]{linear_vb_gibbs_tau_e.png}
  \caption{Posterior comparison for the noise precision $\tau_e$ (linear case).
           The variational Gamma approximation tracks the Gibbs reference closely.}
  \label{fig:linear-posterior-tau}
\end{figure}

\begin{figure}[H]
  \centering
  \includegraphics[width=0.85\textwidth]{linear_mean_errors.png}
  \caption{Posterior mean errors relative to Gibbs reference for all four methods
           (linear case). CAVI and BFGS achieve the smallest errors.}
  \label{fig:linear-errors}
\end{figure}


%─────────────────────────────────────────────────────────────────────────────
\section{Case Study 2: Bayesian Quadratic Regression}
\label{sec:quadratic}
%─────────────────────────────────────────────────────────────────────────────

\subsection{Model}

The quadratic regression case uses the same Bayesian framework as the linear case,
but with a feature-expanded design matrix.
The raw covariate $x_i \sim \mathrm{Uniform}(-2,2)$ is mapped to
$\tilde{x}_i = (1,\; x_i,\; x_i^2)^T$, so the model is
\begin{equation}
  y_i = \beta_0 + \beta_1 x_i + \beta_2 x_i^2 + \varepsilon_i,
  \quad \varepsilon_i \sim \N(0, \tau_e^{-1}).
\end{equation}
The model is linear in $\bm{\beta} = (\beta_0, \beta_1, \beta_2)^T$, so the same
Normal--Gamma ELBO and CAVI updates apply with $p=3$.

The test case uses $n=100$ observations and true parameters
$\bm{\beta}_{\mathrm{true}} = (1.0,\; 2.0,\; -1.0)^T$.
The quadratic term $\beta_2 = -1.0$ introduces a downward curvature that tests whether
the optimisers correctly attribute variance to the additional parameter.

\subsection{ELBO}

The ELBO is identical in form to the linear case (Section~\ref{sec:linear}) with
$p=3$ and the expanded design matrix $\tilde{\bm{X}}$.
There is no new approximation: the feature expansion is a deterministic preprocessing
step applied before inference.

\subsection{Results}

\begin{figure}[H]
  \centering
  \includegraphics[width=0.85\textwidth]{quadratic_data_scatter.png}
  \caption{Synthetic data for Case Study~2: $n=100$ observations, feature expansion
           $(1, x, x^2)$, true $\bm{\beta}=(1.0,2.0,-1.0)$.}
  \label{fig:quadratic-data}
\end{figure}

\begin{figure}[H]
  \centering
  \includegraphics[width=0.85\textwidth]{quadratic_elbo_convergence.png}
  \caption{ELBO convergence for all four optimisation methods on the quadratic test case.
           CAVI again converges fastest; gradient ascent and BFGS follow closely.}
  \label{fig:quadratic-elbo}
\end{figure}

\begin{figure}[H]
  \centering
  \includegraphics[width=0.85\textwidth]{quadratic_vb_gibbs_beta0.png}
  \caption{Posterior comparison for $\beta_0$ (quadratic case).
           Moderate variance collapse: the variational approximation is more
           overconfident than in the linear case.}
  \label{fig:quadratic-posterior-beta0}
\end{figure}

\begin{figure}[H]
  \centering
  \includegraphics[width=0.85\textwidth]{quadratic_vb_gibbs_tau_e.png}
  \caption{Posterior comparison for $\tau_e$ (quadratic case).}
  \label{fig:quadratic-posterior-tau}
\end{figure}

\begin{figure}[H]
  \centering
  \includegraphics[width=0.85\textwidth]{quadratic_mean_errors.png}
  \caption{Posterior mean errors relative to Gibbs reference (quadratic case).
           All gradient-based methods achieve comparable accuracy to CAVI
           on the larger $p=3$ problem.}
  \label{fig:quadratic-errors}
\end{figure}


%─────────────────────────────────────────────────────────────────────────────
\section{Case Study 3: Bayesian Logistic Regression}
\label{sec:logistic}
%─────────────────────────────────────────────────────────────────────────────

\subsection{Model}

The Bayesian logistic regression model is
\begin{align}
  y_i \mid \bm{\beta} &\sim \mathrm{Bernoulli}(\sigma(x_i^T\bm{\beta})),
  \quad i = 1,\ldots,n, \\
  \bm{\beta} &\sim \N(\bm{0},\; \lambda^{-1}\bm{I}_p),
  \quad \lambda^{-1} = 10,
\end{align}
where $\sigma(\psi) = (1+e^{-\psi})^{-1}$ is the logistic sigmoid.
The logistic likelihood is not conjugate to the Gaussian prior, so the posterior is not
available in closed form.

The test case uses $n=200$ observations with $x_i \sim \mathrm{Uniform}(-2,2)$ and
true parameters $\bm{\beta}_{\mathrm{true}} = (-1.0,\; 2.0)^T$ ($p=2$).

\subsection{Jaakkola--Jordan ELBO}

The Jaakkola--Jordan bound \parencite{Jaakkola2000} introduces local variational
parameters $\xi_i > 0$ to replace the intractable log-likelihood term
$\E_q[\log\sigma(x_i^T\bm{\beta})]$ with a tight quadratic lower bound.
At the optimal tightness $\xi_i^* = \sqrt{\E_q[(x_i^T\bm{\beta})^2]}$, the ELBO is
\begin{equation}
  \mathcal{L}(\bm{\theta})
  = \sum_{i=1}^{n}
    \Bigl[(y_i - \tfrac{1}{2})\,x_i^T\bm{\mu}
          - \log 2 - \log\cosh\!\tfrac{\xi_i^*}{2}\Bigr]
  - \mathrm{KL}(q\|p).
\end{equation}
The gradient and CAVI updates for this ELBO are derived in \textcite{Jaakkola2000};
the implementation uses Cholesky reparameterisation with 5 unconstrained parameters.

\subsection{Results}

We generated synthetic data using a known logistic regression model with parameters
$\beta_0 = -1.0$ and $\beta_1 = 2.0$.
These are the data-generating parameters --- they define the process that produced the
observations and represent what we are trying to recover.

Inference was conducted as if the true values were unknown.
The Gibbs sampler uses only two inputs: the 200 binary observations and the weakly
informative prior $\bm{\beta} \sim \N(\bm{0}, 10\bm{I})$.
It never sees the true values.

The posterior distribution is the result of combining what the data say (the likelihood)
with what we believed before seeing the data (the prior).
The posterior mean sits at approximately $(-1.6,\; 2.6)$, offset from the true values,
because with only 200 observations the likelihood is not strong enough to overcome the
prior completely.

The true values appear on the plots only as a benchmark --- to verify that the posterior
is plausible.
The fact that the true values fall within the posterior distribution confirms the sampler
is behaving correctly.

\begin{figure}[H]
  \centering
  \includegraphics[width=0.85\textwidth]{logistic_elbo_convergence.png}
  \caption{ELBO convergence for all four optimisation methods on the logistic test case.
           All methods converge to the same final value within 50 iterations.}
  \label{fig:logistic-elbo}
\end{figure}

\begin{figure}[H]
  \centering
  \includegraphics[width=0.85\textwidth]{logistic_vb_gibbs_beta0.png}
  \caption{Posterior comparison for $\beta_0$: variational approximation versus PG Gibbs
           reference density.
           The posterior mean is offset from the true value $\beta_0=-1.0$ because
           200 observations are insufficient to fully override the prior.}
  \label{fig:logistic-posterior}
\end{figure}

\begin{figure}[H]
  \centering
  \includegraphics[width=0.85\textwidth]{logistic_gibbs_ess.png}
  \caption{Effective sample size (ESS) per parameter for the PG Gibbs sampler.
           ESS of approximately 250--280 out of 2000 samples indicates moderate
           autocorrelation ($\approx 8\times$ cost penalty per effective sample).}
  \label{fig:logistic-ess}
\end{figure}

%─────────────────────────────────────────────────────────────────────────────
\section{Case Study 4: Hierarchical Bayesian Linear Regression}
\label{sec:hierarchical}
%─────────────────────────────────────────────────────────────────────────────

\subsection{Model}

The hierarchical model extends Bayesian linear regression by adding group-level
random effects.  With $J=5$ groups and $N=150$ observations ($\approx 30$ per group):
\begin{align}
  y_{ij} &= x_{ij}^T\bm{\beta} + u_j + \varepsilon_{ij}, \quad
  \varepsilon_{ij} \sim \N(0, \tau_e^{-1}), \\
  u_j &\sim \N(0, \tau_u^{-1}), \qquad j = 1,\ldots,J, \\
  \bm{\beta} &\sim \N(\bm{0}, \lambda^{-1}\bm{I}_p), \quad \lambda = 0.1, \\
  \tau_e &\sim \mathrm{Gamma}(\alpha_e, \gamma_e), \quad
  \tau_u \sim \mathrm{Gamma}(\alpha_u, \gamma_u).
\end{align}
The test case uses $N=150$, $J=5$, $p=2$, $\bm{\beta}_{\mathrm{true}}=(1.0,2.0)^T$,
$\tau_e^{\mathrm{true}}=4.0$, and $\tau_u^{\mathrm{true}}=2.0$.

\subsection{ELBO}

The mean-field approximation factorises as
$q = q(\bm{\beta})\prod_j q(u_j)\,q(\tau_e)\,q(\tau_u)$, giving a closed-form ELBO
with $D=19$ unconstrained parameters: 2 for $\bm{\mu}_\beta$, 3 for the Cholesky factor
of $\bm{\Sigma}_\beta$, 5 group means $m_j$, 5 log-variances $\log v_j$, and 4
log-shape/rate parameters for $q(\tau_e)$ and $q(\tau_u)$.

CAVI updates are all closed form under the conjugate Normal--Gamma structure:
\begin{align}
  q(\bm{\beta}) &= \N\!\left(\bm{\mu}_\beta, \bm{\Sigma}_\beta\right), &
  \bm{\Sigma}_\beta &= \left(\E[\tau_e]\bm{X}^T\bm{X} + \lambda\bm{I}\right)^{-1},\\
  q(u_j) &= \N\!(m_j, v_j), &
  v_j &= \left(\E[\tau_e]n_j + \E[\tau_u]\right)^{-1},\\
  q(\tau_e) &= \mathrm{Gamma}(a_e, b_e), &
  a_e &= \alpha_e + N/2,\\
  q(\tau_u) &= \mathrm{Gamma}(a_u, b_u), &
  a_u &= \alpha_u + J/2.
\end{align}
Full derivations follow the pattern of the linear case (Section~\ref{sec:linear}).

\subsection{Results}

\begin{figure}[H]
  \centering
  \includegraphics[width=0.85\textwidth]{hierarchical_linear_data_scatter.png}
  \caption{Synthetic data for Case Study~4: $N=150$ observations in $J=5$ groups.
           Group membership is colour-coded; the slope and intercept are shared
           ($\bm{\beta}$) while group offsets ($u_j$) vary.}
  \label{fig:hier-data}
\end{figure}

\begin{figure}[H]
  \centering
  \includegraphics[width=0.85\textwidth]{hierarchical_linear_elbo_convergence.png}
  \caption{ELBO convergence for all four optimisation methods (hierarchical case).
           CAVI converges in approximately 100 iterations.
           Newton's method runs to its iteration limit but does not reach the
           CAVI optimum due to the $D=19$ finite-difference Hessian cost.}
  \label{fig:hier-elbo}
\end{figure}

\begin{figure}[H]
  \centering
  \includegraphics[width=0.85\textwidth]{hierarchical_linear_vb_gibbs_beta0.png}
  \caption{Posterior comparison for $\beta_0$ (hierarchical case):
           variational approximation versus three-chain Gibbs reference.
           Both concentrate near the true value of $1.0$.}
  \label{fig:hier-beta0}
\end{figure}

\begin{figure}[H]
  \centering
  \includegraphics[width=0.85\textwidth]{hierarchical_linear_vb_gibbs_tau_u.png}
  \caption{Posterior comparison for the group-precision $\tau_u$ (hierarchical case).
           The variational approximation slightly underestimates posterior spread
           relative to the Gibbs reference.}
  \label{fig:hier-tau-u}
\end{figure}

\begin{figure}[H]
  \centering
  \includegraphics[width=0.85\textwidth]{hierarchical_linear_random_effects.png}
  \caption{Recovered group random effects $\hat{u}_j$ versus true values.
           CAVI and BFGS closely track the Gibbs estimates; Newton's method
           shows larger deviations due to incomplete convergence.}
  \label{fig:hier-re}
\end{figure}

\begin{figure}[H]
  \centering
  \includegraphics[width=0.85\textwidth]{hierarchical_linear_mean_errors.png}
  \caption{Posterior mean errors relative to Gibbs reference (hierarchical case).
           CAVI and BFGS are accurate on $\bm{\beta}$ and $\tau_e$;
           $\tau_u$ is harder to estimate due to only $J=5$ groups.}
  \label{fig:hier-errors}
\end{figure}

%─────────────────────────────────────────────────────────────────────────────
\section{Case Study 5: Hierarchical Bayesian Logistic Regression}
\label{sec:hierarchical-logistic}
%─────────────────────────────────────────────────────────────────────────────

\subsection{Model}

Case Study~5 combines the Jaakkola--Jordan variational framework from Case Study~3
with the hierarchical random-effects structure from Case Study~4.
With $J=5$ groups and $N=150$ observations:
\begin{align}
  y_{ij} &\sim \mathrm{Bernoulli}\!\left(\sigma(x_{ij}^T\bm{\beta} + u_j)\right), \\
  u_j &\sim \N(0, \tau_u^{-1}), \qquad j = 1,\ldots,J, \\
  \bm{\beta} &\sim \N(\bm{0}, \lambda^{-1}\bm{I}_p), \quad \lambda = 0.1, \\
  \tau_u &\sim \mathrm{Gamma}(\alpha_u, \gamma_u).
\end{align}
The test case uses $N=150$, $J=5$, $p=2$, $\bm{\beta}_{\mathrm{true}}=(-1.0,2.0)^T$,
and $\tau_u^{\mathrm{true}}=2.0$.

\subsection{ELBO}

The mean-field approximation factorises as
$q = q(\bm{\beta})\prod_j q(u_j)\,q(\tau_u)$.
The Jaakkola--Jordan bound introduces variational parameters $\xi_{ij}$ per
observation to restore tractability:
\begin{equation}
  \mathcal{L} \approx \sum_{ij} \E_q[\log \sigma_{\xi_{ij}}(x_{ij}^T\bm{\beta}+u_j)]
    - \mathrm{KL}(q(\bm{\beta})\|p(\bm{\beta}))
    - \sum_j \mathrm{KL}(q(u_j)\|p(u_j))
    - \mathrm{KL}(q(\tau_u)\|p(\tau_u)),
\end{equation}
where $\sigma_\xi$ denotes the JJ surrogate likelihood.
The ELBO has $D=17$ unconstrained parameters: 2 for $\bm{\mu}_\beta$, 3 for the
Cholesky factor of $\bm{\Sigma}_\beta$, 5 group means $m_j$, 5 log-variances
$\log v_j$, and 2 log-shape/rate parameters for $q(\tau_u)$.

CAVI updates follow the JJ fixed-point equations:
\begin{align}
  \xi_{ij}^2 &= \E_q[(x_{ij}^T\bm{\beta}+u_j)^2], \\
  q(\bm{\beta}) &= \N(\bm{\mu}_\beta, \bm{\Sigma}_\beta), \quad
    \bm{\Sigma}_\beta = \left(2\sum_{ij}\lambda(\xi_{ij})x_{ij}x_{ij}^T + \lambda\bm{I}\right)^{-1}, \\
  q(u_j) &= \N(m_j, v_j), \quad
    v_j = \left(2\sum_i \lambda(\xi_{ij}) + \E[\tau_u]\right)^{-1},
\end{align}
where $\lambda(\xi) = (\sigma(\xi)-\tfrac{1}{2})/(2\xi)$ is the JJ function.

\subsection{Results}

\begin{figure}[H]
  \centering
  \includegraphics[width=0.85\textwidth]{hierarchical_logistic_data_scatter.png}
  \caption{Synthetic data for Case Study~5: $N=150$ binary observations in $J=5$ groups.
           Group membership is colour-coded; positive outcomes shown as filled circles.}
  \label{fig:hier-log-data}
\end{figure}

\begin{figure}[H]
  \centering
  \includegraphics[width=0.85\textwidth]{hierarchical_logistic_elbo_convergence.png}
  \caption{ELBO convergence for all four optimisation methods (hierarchical logistic case).
           CAVI converges in approximately 37 iterations.
           Newton's method fails to reach the CAVI optimum
           (ELBO $= -77.4$ vs\ $-54.9$ for CAVI/BFGS).}
  \label{fig:hier-log-elbo}
\end{figure}

\begin{figure}[H]
  \centering
  \includegraphics[width=0.85\textwidth]{hierarchical_logistic_vb_gibbs_beta0.png}
  \caption{Posterior comparison for $\beta_0$ (hierarchical logistic case):
           variational approximation versus P\'{o}lya--Gamma Gibbs reference.
           Both concentrate near the true value of $-1.0$.}
  \label{fig:hier-log-beta0}
\end{figure}

\begin{figure}[H]
  \centering
  \includegraphics[width=0.85\textwidth]{hierarchical_logistic_vb_gibbs_beta1.png}
  \caption{Posterior comparison for $\beta_1$ (hierarchical logistic case):
           variational approximation versus Gibbs reference near the true value of $2.0$.}
  \label{fig:hier-log-beta1}
\end{figure}

\begin{figure}[H]
  \centering
  \includegraphics[width=0.85\textwidth]{hierarchical_logistic_sd_ratios.png}
  \caption{Posterior standard deviation ratios (VB/Gibbs) for the hierarchical
           logistic case. Both $\beta_0$ and $\beta_1$ show symmetric collapse
           (SD ratio $\approx 0.46$), in contrast to the asymmetric collapse
           observed in Case Study~4.}
  \label{fig:hier-log-sd}
\end{figure}

\begin{figure}[H]
  \centering
  \includegraphics[width=0.85\textwidth]{hierarchical_logistic_timing_comparison.png}
  \caption{Timing comparison for the hierarchical logistic case.
           CAVI is 1050$\times$ faster than PG Gibbs; BFGS is 81$\times$ faster.
           The PG Gibbs sampler requires $n\times T$ Gamma draws per iteration
           ($N=150$, 3$\times$2000 samples).}
  \label{fig:hier-log-timing}
\end{figure}

%─────────────────────────────────────────────────────────────────────────────
\section{Discussion}
\label{sec:discussion}
%─────────────────────────────────────────────────────────────────────────────

\subsection{Cross-Case Comparison}

Table~\ref{tab:comparison} summarises the key structural and empirical differences
across the five test cases.
All timing figures are from the same machine (Alienware 17~R5, Intel i7-8750H) but
are not controlled comparisons: model complexity, parameter count, and Gibbs overhead
differ across cases.
The fair within-notebook comparison is each VI method versus its own Gibbs reference.

\begin{table}[H]
\centering
\caption{Comparison of the five case studies.}
\label{tab:comparison}
\begin{tabular}{@{}llllll@{}}
\toprule
Property & Linear & Quadratic & Logistic & Hier.~Linear & Hier.~Logistic \\
\midrule
Likelihood & Gaussian & Gaussian & Bernoulli & Gaussian & Bernoulli \\
True $\bm{\beta}$ & $(1.0, 2.0)$ & $(1.0, 2.0, -1.0)$ & $(-1.0, 2.0)$ & $(1.0, 2.0)$ & $(-1.0, 2.0)$ \\
Predictors $p$ & 2 & 3 & 2 & 2 & 2 \\
Observations $n$ & 50 & 100 & 200 & 150 ($J\!=\!5$) & 150 ($J\!=\!5$) \\
Conjugacy & Yes & Yes & No (JJ) & Yes & No (JJ) \\
ELBO form & Normal--Gamma & Normal--Gamma & JJ bound & Normal--Gamma & JJ + hier.\ \\
Reference sampler & Conj.~Gibbs & Conj.~Gibbs & PG Gibbs & Hier.~Gibbs & PG Gibbs \\
Gibbs cost & Low & Low & High & Medium & High \\
Vartnl.~params $D$ & 8 & 11 & 5 & 19 & 17 \\
Variance collapse & Mild & Moderate & Minimal & Mild ($\beta_0$) & Symmetric \\
CAVI updates & Closed form & Closed form & Via JJ & Closed form & Via JJ \\
\bottomrule
\end{tabular}
\end{table}

\subsection{Posterior Variance Collapse}

Posterior variance collapse --- the tendency of mean-field VI to underestimate posterior
uncertainty --- varies systematically across the five cases.

In the linear and quadratic Gaussian cases, the mean-field factorisation
$q(\bm{\beta})\,q(\tau_e)$ removes true posterior correlations between $\bm{\beta}$ and
$\tau_e$, squeezing the marginal approximations.
The effect is mild for the linear case (small $p=2$ and small $n=50$) and more pronounced
for the quadratic case (larger $p=3$ and $n=100$).

In the logistic case with the Jaakkola--Jordan bound, variance collapse is minimal.
This is because the JJ bound is tight at the CAVI fixed point: when
$\xi_i^* = \sqrt{\E_q[(x_i^T\bm{\beta})^2]}$, the bound equals the true expectation,
and the ELBO objective penalises over-concentration.
The result is that CAVI and all gradient-based methods recover posterior standard
deviations accurately relative to the PG Gibbs reference.

In the hierarchical linear case, the mean-field factorisation separates $q(\bm{\beta})$
from $q(u_j)$ and both from $q(\tau_u)$, removing correlations between the fixed
effects and the random-effect precision.
The resulting collapse is asymmetric: severe for $\beta_0$ (SD ratio $0.130$) but
negligible for $\beta_1$ (SD ratio $0.987$), because the intercept is entangled
with the group random effects while the slope is not.

In the hierarchical logistic case, variance collapse is symmetric across both
coefficients (SD ratio $\approx 0.46$ for both $\beta_0$ and $\beta_1$).
The JJ bound removes the asymmetry seen in Case Study~4: without a Gaussian precision
parameter $\tau_e$, the mean-field factorisation treats $\bm{\beta}$ and $u_j$ more
evenly, producing balanced underestimation.

\subsection{Optimisation Behaviour}

CAVI is the most efficient method across all five settings, because it exploits
closed-form coordinate updates.
For the linear, quadratic, and hierarchical linear cases, the conjugate Normal--Gamma
structure enables exact updates; for the logistic and hierarchical logistic cases,
the JJ bound restores the same conjugate structure.

BFGS converges rapidly in all five cases, providing an effective gradient-based
alternative when closed-form CAVI updates are not available.

Newton's method is the most expensive per-iteration due to the finite-difference
Hessian.
For the low-dimensional cases ($D \leq 11$) it converges cleanly; for the hierarchical
cases ($D=19$ and $D=17$), the Hessian cost is higher and the method fails to reach
the CAVI optimum, making BFGS the preferred second-order alternative.

%─────────────────────────────────────────────────────────────────────────────
\section*{Conclusion}
\addcontentsline{toc}{section}{Conclusion}
%─────────────────────────────────────────────────────────────────────────────

This project has applied gradient-based optimisation to variational inference across
five Bayesian regression models of increasing difficulty, using ELBO maximisation to
turn intractable posterior inference into a problem directly amenable to the methods
covered in ENEL445.

The five cases form a natural progression.
The linear case establishes the baseline: a conjugate Gibbs reference, mild variance
collapse, and clean convergence for all four methods.
Moving to quadratic regression ($p=3$) raises the parameter count and shows that
moderate collapse can persist even after the ELBO has converged.
The logistic case introduces a non-conjugate bottleneck, resolved by the
Jaakkola--Jordan bound, and reveals that collapse can be avoided almost entirely when
that bound is tight at the CAVI fixed point.
Adding group random effects in the hierarchical linear case ($J=5$, $D=19$) shows the
same framework scaling to mixed-effects models; the intercept collapse is severe and
asymmetric, traced back to the entanglement of $\beta_0$ with the group offsets.
The hierarchical logistic case ($D=17$) closes the loop by combining the JJ bound
with random effects: collapse becomes symmetric across both coefficients, and PG Gibbs
confirms the variational solution.

Across all five settings, CAVI is the most efficient optimiser, exploiting closed-form
coordinate updates wherever conjugacy permits.
BFGS is the natural gradient-based fallback --- it reaches near-identical solutions
using only gradient evaluations and Wolfe-condition step control, without requiring
closed-form updates.

The P\'{o}lya--Gamma Gibbs sampler provides accurate reference posteriors for the
logistic and hierarchical logistic cases, but it is expensive: $n\times T$ Gamma draws
per iteration, compared to CAVI's closed-form sweeps.
For settings where $n$ is large, the computational gap is substantial --- CAVI is
1050$\times$ faster than PG Gibbs in Case Study~5 --- which gives VI methods a
practical advantage even when exact inference is in principle available.

\nocite{*}
\printbibliography

%─────────────────────────────────────────────────────────────────────────────
\appendix
%─────────────────────────────────────────────────────────────────────────────

\section{Case Study Summary Table}
\label{app:comparison-table}

The following table expands on Table~\ref{tab:comparison} with additional
optimisation-specific and sampler-specific properties.

\begin{table}[H]
\centering
\caption{Extended case study comparison.}
\begin{tabular}{@{}llllll@{}}
\toprule
Property & Linear & Quadratic & Logistic & Hier.~Lin. & Hier.~Log. \\
\midrule
\multicolumn{6}{@{}l}{\textbf{Model}} \\
Likelihood & Gaussian & Gaussian & Bernoulli & Gaussian & Bernoulli \\
True $\bm{\beta}$ & $(1.0, 2.0)$ & $(1.0, 2.0, -1.0)$ & $(-1.0, 2.0)$ & $(1.0, 2.0)$ & $(-1.0, 2.0)$ \\
Predictors $p$ & 2 & 3 & 2 & 2 & 2 \\
Observations $n$ & 50 & 100 & 200 & 150 ($J\!=\!5$) & 150 ($J\!=\!5$) \\
Conjugacy & Yes & Yes & No & Yes & No \\
\midrule
\multicolumn{6}{@{}l}{\textbf{Variational approximation}} \\
Variational family & Normal--Gamma & Normal--Gamma & Gaussian & Normal--Gamma & Gaussian \\
ELBO form & Closed form & Closed form & JJ bound & Closed form & JJ + hier.\ \\
VI parameters $D$ & 8 & 11 & 5 & 19 & 17 \\
Variance collapse & Mild & Moderate & Minimal & Asymmetric & Symmetric \\
\midrule
\multicolumn{6}{@{}l}{\textbf{Reference sampler}} \\
Sampler type & Conj.~Gibbs & Conj.~Gibbs & PG Gibbs & Hier.~Gibbs & PG Gibbs \\
Cost per sweep & $O(p^3)$ & $O(p^3)$ & $O(n \cdot T)$ & $O(N + J)$ & $O(n \cdot T)$ \\
ESS (typical) & High & High & $\approx 250$--$280$ & High & Moderate \\
\midrule
\multicolumn{6}{@{}l}{\textbf{Optimisation}} \\
CAVI convergence & $<50$ iter & $<50$ iter & $<50$ iter & $\approx 100$ iter & $\approx 37$ iter \\
BFGS convergence & Fast & Fast & Fast & Fast & Fast \\
Recommended method & CAVI & CAVI & CAVI & CAVI & CAVI \\
\bottomrule
\end{tabular}
\end{table}

The text associated with the logistic case for presentation

\begin{framed}
\noindent
We generated synthetic data using a known logistic regression model with parameters
$\beta_0 = -1.0$ and $\beta_1 = 2.0$.
These are the data-generating parameters: they define the process that produced the
observations and represent what we are trying to recover.

We then treated inference as if the true values were unknown.
The Gibbs sampler uses only two inputs: the 200 binary observations and a weakly
informative prior $\bm{\beta} \sim \N(\bm{0}, 10\bm{I})$.
It never sees the true values.

The posterior distribution is the result of combining what the data say (the likelihood)
with what we believed before seeing the data (the prior).
The posterior mean sits at approximately $(-1.6, 2.6)$, offset from the true values,
because with only 200 observations the likelihood is not strong enough to overcome
the prior completely.

The true values appear on the plots only as a benchmark --- to verify that the posterior
is plausible.
The fact that the true values fall within the posterior distribution confirms the sampler
is behaving correctly.
\end{framed}

%─────────────────────────────────────────────────────────────────────────────
\section{Individual Report References}
\label{app:individual-refs}
%─────────────────────────────────────────────────────────────────────────────

The five individual case study reports contain full derivations, all figures, and
complete code listings.

\begin{description}
  \item[CS1 --- Linear regression:] \texttt{report/20260514A-LINEAR-CASE-STUDY.tex}.
    Full ELBO derivation, gradient verification, ELBO landscape, all convergence
    figures, posterior comparison figures, and variance collapse analysis.
  \item[CS2 --- Quadratic regression:] \texttt{report/20260514A-QUADRATIC-CASE-STUDY.tex}.
    Identical framework to the linear case with $p=3$ and feature expansion.
  \item[CS3 --- Logistic regression:] \texttt{report/20260514A-LOGISTIC-CASE-STUDY.tex}.
    Full Jaakkola--Jordan derivation, P\'{o}lya--Gamma sampler, ESS diagnostics,
    and all 20 result figures.
  \item[CS4 --- Hierarchical linear regression:] \texttt{report/20260514A-HIERARCHICAL-LINEAR-CASE-STUDY.tex}.
    Two-level Normal--Gamma ELBO, group random effects, three-chain Gibbs sampler,
    and all 23 result figures. Asymmetric $\beta_0$ variance collapse.
  \item[CS5 --- Hierarchical logistic regression:] \texttt{report/20260514A-HIERARCHICAL-LOGISTIC-CASE-STUDY.tex}.
    JJ bound with random effects, P\'{o}lya--Gamma Gibbs reference,
    and all 20 result figures. Symmetric variance collapse across both coefficients.
\end{description}

%─────────────────────────────────────────────────────────────────────────────
\section{Code Structure}
\label{app:code}
%─────────────────────────────────────────────────────────────────────────────

All code is in the \texttt{python/} directory.
Each notebook follows a five-stage structure.

\begin{description}
  \item[Stage 1:] Data generation. Saves \texttt{results/data/<case>\_data.parquet}.
  \item[Stage 2:] ELBO evaluation and gradient verification (finite-difference check).
  \item[Stage 3:] Optimisation (CAVI, gradient ascent, Newton, BFGS).
    Saves \texttt{results/data/<case>\_opt.parquet}.
  \item[Stage 4:] Reference Gibbs sampler with ESS diagnostics.
    Saves \texttt{results/data/<case>\_gibbs.parquet}.
  \item[Stage 5:] Diagnostics and comparison figures.
    Saves \texttt{results/tables/<case>\_diagnostics.csv}.
\end{description}

\begin{center}
\begin{tabular}{@{}lll@{}}
\toprule
Case & Notebook & Figure prefix \\
\midrule
Linear          & \texttt{python/linear\_run.ipynb}                  & \texttt{linear\_} \\
Quadratic       & \texttt{python/quadratic\_run.ipynb}               & \texttt{quadratic\_} \\
Logistic        & \texttt{python/logistic\_run.ipynb}                & \texttt{logistic\_} \\
Hier.~Linear    & \texttt{python/hierarchical\_linear\_run.ipynb}    & \texttt{hierarchical\_linear\_} \\
Hier.~Logistic  & \texttt{python/hierarchical\_logistic\_run.ipynb}  & \texttt{hierarchical\_logistic\_} \\
\bottomrule
\end{tabular}
\end{center}

Core algorithms are in \texttt{python/vb\_algorithms\_py.py};
utilities (archive, plot, I/O) are in \texttt{python/vb\_utils\_py.py}.

All five individual case study reports are available online at:
\begin{center}
  \url{https://ENEL445.IndustrialOrigami.ai}
\end{center}

\end{document}
